\documentclass[11pt]{article}
\usepackage[T1]{fontenc}
\usepackage[utf8]{inputenc}
\usepackage{lmodern,microtype,amsmath,amssymb,amsthm,mathtools}
\usepackage[letterpaper,margin=0.92in,headheight=15pt]{geometry}
\usepackage{booktabs,longtable,array,enumitem,xcolor,fancyhdr}
\usepackage{hyperref}
\definecolor{ink}{HTML}{183650}
\definecolor{accent}{HTML}{16666C}
\hypersetup{colorlinks=true,linkcolor=accent,urlcolor=accent,citecolor=accent,
pdftitle={Smaller explicit Diophantine equations for ZFC consistency},
pdfauthor={AI-assisted research prepared for Vladimir Reshetnikov}}
\pagestyle{fancy}\fancyhf{}
\fancyhead[L]{\small Smaller Diophantine consistency equations}
\fancyhead[R]{\small ProveIt research}
\fancyfoot[C]{\thepage}
\setlength{\parindent}{0pt}\setlength{\parskip}{0.4em}
\setlength{\emergencystretch}{3em}
\setlist{nosep,leftmargin=1.5em}
\allowdisplaybreaks[2]
\newcommand{\N}{\mathbb N}
\newcommand{\Z}{\mathbb Z}
\newcommand{\Pos}{\mathbb Z_{>0}}
\newcommand{\Con}{\operatorname{Con}}
\newcommand{\pc}{\operatorname{pc}}
\newcommand{\Sq}{\operatorname{Sq}}
\newcommand{\file}[1]{\texttt{\small\detokenize{#1}}}
\newtheorem{theorem}{Theorem}[section]
\newtheorem{lemma}[theorem]{Lemma}
\newtheorem{proposition}[theorem]{Proposition}
\theoremstyle{definition}\newtheorem{definition}[theorem]{Definition}
\theoremstyle{remark}\newtheorem{remark}[theorem]{Remark}
\newcommand{\NineBestDegree}{15\,534\,796}
\newcommand{\NineBestOps}{655}
\newcommand{\TenBestDegree}{39\,914\,856}
\newcommand{\TenBestOps}{651}

\title{\color{ink}Smaller explicit Diophantine equations\\for ZFC consistency\\[0.7em]
\large Encoding the monomials that occur, overlapping output positions,\\and polynomial approximations to square roots}
\author{AI-assisted mathematical research\\Prepared for Vladimir Reshetnikov and the ProveIt project}
\date{23 October 2026}
\begin{document}
\maketitle
\begin{abstract}
We construct explicit integer polynomials whose unsolvability represents
the syntactic consistency of ZFC, improving the current ProveIt constructions
on both degree and arithmetic cost. The main nine-positive-variable equation
has exact degree \NineBestDegree{} and a division-free circuit with
\NineBestOps{} live binary operations; its nine-nonnegative-variable version
uses three more operations. A ten-unrestricted-integer equation has exact
degree \TenBestDegree{} and \TenBestOps{} operations. Lower-cost alternatives
are supplied. The comparison is with the latest wave-6 constructions in
repository commit \texttt{124c07c0b77a}, rather than a superseded baseline.

The first new observation is that a quadratic source need not distinguish
all possible monomials: it must distinguish each monomial that occurs from
all competing monomials. For the actual source this reduces the largest
position from 799 to 179. A second finite certificate permits overlapping
row blocks. A third improvement uses divisibility already present in the
Jones arithmetic witnesses to construct a closer polynomial majorant for
the radical in the final norm. All new arithmetic arguments are proved;
the archive contains the literal circuits, complete fixed coefficient data,
independent checks, and the inherited theorem dependencies. The resulting
consistency equivalence is an ordinary mathematical proof formalizable in
ZFC, using the same native compiler and machine calibration as the supplied
repository. We do not claim a new end-to-end Lean/Rocq certificate, a proof of
consistency, or a global minimum over Diophantine equations.
\end{abstract}
\newpage
\tableofcontents
\newpage

\section{The result, its domain, and the comparison}
\label{sec:result}
Write $\N=\{0,1,2,\ldots\}$ and $\Pos=\{1,2,\ldots\}$.
We use the repository's precise syntactic consistency sentence
\[
 \Con(\mathrm{ZFC})\;\equiv\;
 \neg\exists p\;\operatorname{Proof}_{\mathrm{ZFC}}(p,\bot),
\]
calibrated there by \texttt{BProv ZFCax\_s [] Form.fBot}.
The coefficient tuple is fixed independently of this sentence's truth.

\begin{theorem}[Explicit consistency equations]
The finite data and polynomial formulas in this article define polynomials
$F_9\in\Z[x_1,\ldots,x_9]$ and $F_{10}\in\Z[z_1,\ldots,z_{10}]$ such that,
using the inherited mathematical compiler and arithmetic theorems identified
in Section~\ref{sec:foundations}, ZFC proves
\begin{align*}
 \neg\exists\mathbf x\in\Pos^9\;F_9(\mathbf x)=0
 &\ \longleftrightarrow\ \Con(\mathrm{ZFC}),\\
 \neg\exists\mathbf z\in\Z^{10}\;F_{10}(\mathbf z)=0
 &\ \longleftrightarrow\ \Con(\mathrm{ZFC}).
\end{align*}
The principal $F_9$ has exact total degree \NineBestDegree{} and
\NineBestOps{} binary operations. The principal $F_{10}$ has exact total
degree \TenBestDegree{} and \TenBestOps{} operations. The displayed
registers in their definitions are abbreviations, not additional unknowns.
\end{theorem}

\begin{table}[htbp]
\centering\small
\begin{tabular}{@{}llrr@{}}
\toprule Construction & Domain & Exact degree & Operations\\\midrule
Package 45 (previous) & $\Pos^9$ & $129\,875\,980$ & 788\\
New source, uniform, inherited tail & $\Pos^9$ & $29\,157\,244$ & 562\\
New source, uniform, V8 & $\Pos^9$ & $24\,399\,300$ & 568\\
New source, uniform, principal tail & $\Pos^9$ & $20\,754\,916$ & 599\\
New source, overlap, V8 & $\Pos^9$ & $18\,262\,476$ & 624\\
\textbf{New source, overlap, principal} & $\boldsymbol{\Pos^9}$ & $\mathbf{15\,534\,796}$ & \textbf{655}\\\midrule
Package 40 (previous) & $\Z^{10}$ & $237\,520\,456$ & 782\\
New source, uniform & $\Z^{10}$ & $53\,325\,896$ & 595\\
\textbf{New source, overlap} & $\boldsymbol{\Z^{10}}$ & $\mathbf{39\,914\,856}$ & \textbf{651}\\
\bottomrule
\end{tabular}
\caption{The new degree--operation tradeoffs. Every nine-variable natural-domain
version has the same degree and three additional operations. The archive also
includes the overlapping source with the inherited tail (618 operations,
degree $21\,823\,612$). Certificate-circuit operations are not charged to
the primary equation.}
\end{table}

Relative to the latest comparator, the principal nine-variable degree is
smaller by a factor of approximately $8.3603$, and its operation count falls
by 133 ($16.88\%$). The principal ten-integer degree is smaller by a factor
of approximately $5.9507$, and its operation count falls by 131 ($16.75\%$).


The nine-variable result retains positivity. Its unrestricted-integer
interpretation is not asserted. The supplied natural-variable version uses
an explicit coordinate shift described below. For a Diophantine equation
over all integers, $F_{10}$ is the corresponding primary deliverable.
That branch retains Sun's ten-integer reduction \cite{sun}; the
nine-variable branch uses the older Jones arithmetic bridge \cite{jones}.

The preceding nine-positive-variable point in package 45 is degree
$129\,875\,980$ at 788 operations. The preceding ten-integer point in
package 40 is degree $237\,520\,456$ at 782 operations
\cite{package40,package45}. We compare against those latest points. Earlier
claims comparing only with exponentially larger wave-5 polynomials are
not the baseline for this contribution.

There is no metric-independent ``simplest'' equation. The operation count
charges each binary addition, subtraction, and multiplication, including
multiplication by a fixed coefficient. Reuse of a computed register is free;
fixed powers are expanded into charged multiplication gates. The count does
not measure bit complexity, expanded monomial count, or the cost of producing
fixed coefficients. The repository's 67-operation equation with fifteen
positive variables remains cheaper on that single measure. We have not
found a sub-nine-natural-variable or sub-ten-integer construction.

\section{The fixed quadratic parent and coefficient data}
\label{sec:parent}
The common parent has 31 strictly positive coordinates and 22 quadratic
equations. In order, the coordinates are
\[
\begin{gathered}
J,F,\alpha,z,f,h,i,o,j,s,w,\tau,\eta,\zeta,y,\gamma,\\
N,E,P,c,a,X,D,C_2,L,F_2,V,g,I,B_a,\Delta.
\end{gathered}
\]
Here these letters are local to the parent; later tail registers reuse some
letters for different quantities. The four fixed compiler ports are
$(\beta,K_0,M_C,M_F)$, with the high mask $M_F$ already shifted. Put
\[
q=\beta J+1,\qquad k=\eta+\zeta,\qquad
H=4a+3,\qquad C=q-F-\alpha.
\]
The equations, in their original order, are
\begin{longtable}{@{}r@{\quad}l@{}}
\toprule Row & Equation\\\midrule\endhead
1& $L+1=q^2$\\
2& $N=q(L+1)$\\
3& $X=wq$\\
4& $a-E=sN$\\
5& $E=X(a-E)$\\
6& $c-\eta=k(a-E)$\\
7& $P-k=E(c-\eta)$\\
8& $\tau^2=P(P-k)+1$\\
9& $D=X+ac+\gamma H$\\
10& $\Delta=a^2+H$\\
11& $C_2=c^2$\\
12& $D^2=\Delta C_2+1$\\
13& $k=jc-V+1+hE$\\
14& $jc-V=(q-F)(N-q)-CL+(M_C+qM_F)J$\\
15& $(K_0+w)C=F+z\beta J$\\
16& $F_2=f^2$\\
17& $V=of-c$\\
18& $g=y^2-V^2$\\
19& $I=iC_2$\\
20& $B_a=I^2$\\
21& $B_a=\Delta(F_2-1)$\\
22& $y^2-B_ag=1$\\\bottomrule
\end{longtable}

The exact orientation of each residual is fixed by
\file{data/parent_system.json}. The displayed equalities alone are
unchanged by reversing a residual; the finite coefficient file records
the orientation used in the packing. For each parent coordinate substitute
$x_i+1$ and form the recorded left-minus-right residual. Denote these
polynomials by $r_1,\ldots,r_{22}$. Every $r_j$ has degree at most two.
The fifth residual has $r_5(\mathbf0)=1$.

The imported parent theorem is
\begin{equation}\label{eq:parent}
 \exists\mathbf x\in\N^{31}\ \bigwedge_{j=1}^{22}r_j(\mathbf x)=0
 \quad\longleftrightarrow\quad
 \exists p\;\operatorname{Proof}_{\mathrm{ZFC}}(p,\bot).
\end{equation}
We use the complete package-31 tuple, whose radix is $B_0=2^{1064921}$.
Its four port bit lengths are respectively
\[
1064921,\quad1070004,\quad1064921,\quad1064922.
\]
Their compressed unsigned little-endian payloads are included, together
with hashes and the inherited coefficient recipe. They total 533,099
uncompressed bytes. No component is taken from a different compiler tuple.
The archive also materializes all 16 fixed leaves of the emitted graphs
in \file{data/fixed/}, including signs, exact byte lengths, and SHA-256
hashes; these payloads contain 3,331,623 uncompressed bytes in total.

Every additional fixed leaf is either a literal small integer, the explicit
integer $2^{2140030}$, or a recorded integer polynomial in these four ports.
Thus the equation contains no unspecified recursively enumerable index,
consistency-dependent coefficient, or search for a ZFC refutation in its
definition. A gate file plus these finite leaves is a complete specification
of one ordinary integer polynomial, even though fully expanding it would be
an inefficient way to display it.

\section{A smaller positional certificate}
\label{sec:packing}
\subsection{Only the present monomials need isolation}
Adjoin $x_0=1$ and homogenize each monomial of degree at most two to
degree two. There are $32\cdot33/2=528$ possible unordered pairs of indices.
Only 71 of those pairs occur in the 22 residuals. We call these pairs
\emph{present}. Assign the following weights in the coordinate order
$(0,J,F,\ldots,\Delta)$:
\begin{equation}\label{eq:weights}
\begin{split}
(&0,174,166,110,137,115,65,168,126,105,49,155,54,60,15,4,\\
 &118,146,26,44,7,17,57,38,158,177,179,116,46,94,35,96).
\end{split}
\end{equation}
Let $W=179$. All nonpadding weights are distinct and positive.
For every present pair $(a,b)$, its sum $w_a+w_b$ differs from the sum
of every other pair among the 528. Pairs absent from every residual may
collide. Indeed, these 528 pairs occupy just 266 different positions.

The certificate is a finite list and a universally quantified check over
528 pairs, not a claim of optimality for a combinatorial search. Both the
construction verifier and a separately written sparse-polynomial checker
expand the actual parent to identify the 71 present pairs and check this
property. No heuristic search code is needed to validate it.

\subsection{Extraction with possibly overlapping row blocks}
Write $r_j=\sum_\alpha p_{j,\alpha}x^\alpha$, where $|\alpha|\le2$, and set
\[
 s(\alpha)=\sum_i\alpha_iw_i,\qquad
 c_{j,\alpha}=(2-|\alpha|)!\prod_i\alpha_i!\;p_{j,\alpha}.
\]
For a row permutation $\pi$, signs $\sigma_j\in\{1,-1\}$, and increasing
integer targets $H_j\ge2W$, define
\begin{equation}\label{eq:Dcode}
 \mathcal D(T)=\sum_{j=0}^{21}\sum_{\alpha\in\operatorname{supp}r_{\pi(j)}}
 \sigma_jc_{\pi(j),\alpha}T^{H_j-s(\alpha)}.
\end{equation}
Every exponent here is nonnegative and fixed.

\begin{lemma}[Isolation only at the selected outputs]\label{lem:isolation}
Suppose for every row $j$, every $\alpha\in\operatorname{supp}r_{\pi(j)}$,
every degree-at-most-two monomial $\beta$, and every row $i$, the equality
\[
 H_i-H_j=s(\beta)-s(\alpha)
\]
implies $i=j$ and $\beta=\alpha$. Then, for
$\mathcal C(T)=(1+\sum_i x_iT^{w_i})^2$,
\[
 [T^{H_j}]\,\mathcal C(T)\mathcal D(T)
       =2\sigma_jr_{\pi(j)}(\mathbf x).
\]
Moreover all monomial records in $\mathcal D$ have distinct exponents.
\end{lemma}
\begin{proof}
A term from row $j$ and monomial $\alpha$ can contribute at output $H_i$
only under the displayed equality. The isolation hypothesis leaves exactly
the matching monomial in the matching row. Its coefficient in $\mathcal C$
is $2!\,x^\alpha/((2-|\alpha|)!\prod_i\alpha_i!)$, cancelling the
factorial multiplier in $c_{j,\alpha}$. This argument also treats constants
and linear terms by the padding variable. If two records of $\mathcal D$
shared an exponent, multiplying either by its own present monomial from
$\mathcal C$ would violate output isolation at that row. Hence the records
have distinct exponents.
\end{proof}

The condition distinguishes a present monomial from \emph{every} competing
monomial. Distinguishing present monomials only from one another would not
be sufficient: an absent monomial is still produced by $\mathcal C$ and
could contaminate a selected coefficient.

\subsection{Two explicit placements}
The uniform placement uses row order
$1,2,3,4,6,\ldots,22,5$, all signs positive, and
\[
 H_j=358+359j\quad(0\le j\le21),\qquad H_{21}=7897.
\]
Distinct row targets are more than $2W$ apart, so the 71-pair certificate
proves the lemma's hypothesis immediately. The highest coefficient of
$\mathcal D$ is $2r_5(\mathbf0)=2$.

The overlapping placement has targets
\begin{equation}\label{eq:overlap}
\begin{split}
(&358,608,939,1147,1432,1739,1989,2170,2512,2699,2972,\\
 &3305,3492,3683,3914,4108,4399,4697,4926,5224,5469,5775).
\end{split}
\end{equation}
Its original row order and signs are specified in
\file{data/overlap_layout.json}; the finite table is reproduced in the
appendix. The last row is the negative of original row 9 and has constant
$+8$, so the highest coefficient of $\mathcal D$ is $16$.
Every successive gap is greater than $W$. Differences of nonadjacent
targets therefore exceed $2W$, and only adjacent pairs require the finite
cross-row check. That check verifies the full hypothesis of
Lemma~\ref{lem:isolation}. No claim that all support intervals are disjoint
is made or needed.

\section{From isolated coefficients to the arithmetic source}
\label{sec:source}
We spell out the source because replacing positions inside a universal
theorem with hard-coded exponents would be invalid. The new source is
proved directly and then supplied to the numerical arithmetic bridge.

Let $H_*=\max_jH_j$, $K_*=H_*+2W$, and let $M_*$ bound the absolute values
of every weighted coefficient $c_{j,\alpha}$. The finite coefficient check
gives $M_*<2^{2134927}$. The chosen fixed integer
$\Lambda=2^{2140030}$ comfortably satisfies all coefficient bounds below.
For the nine-variable construction use
\[
 b=\Lambda+\epsilon,\quad B=b^3,\quad\epsilon\ge0;
\]
for the ten-integer construction use $b=1+3f$, $B=\Lambda b^2$.
The latter bridge first recovers $b\ge4$ before the digit argument is used.

Define the following polynomial registers, with $g$ the encoded witness:
\begin{align*}
 \ell&=\sum_{i=1}^{31}B^{w_i}, &
 N_0&=B^{W+1}, & N_1&=4B^{K_*+1}, & N&=N_0N_1,\\
 M_0&=N_0-1-(b-1)\ell, &
 T_*&=(B-2)\sum_jB^{H_j},\\
 \mathcal T&=M_0+T_*N_0, &
 \mathcal J&=2(1+g)^2\mathcal D(B)+B\sum_{v=0}^{K_*}B^v,\\
 \mathcal S&=g+\mathcal JN_0, &
 R&=(N-1)\bigl((\mathcal S+\mathcal T+1)N+\mathcal T+1\bigr).
\end{align*}
The source cap is $G_*=32bB^W$ in the nine-variable case and
$C_*=bB^W$ in the ten-integer case. All sums have fixed finite limits.

\begin{lemma}[The carry test survives the sparse encoding]
The source has the same two-way interface as the source in package 45:
if $b$ is a power of two and $0<g<4bB^W$, then
\[
 N^2\mid\binom{2R}{R}
 \quad\Longrightarrow\quad
 \exists\mathbf x\in\N^{31}\ \bigwedge_j r_j(\mathbf x)=0.
\]
Conversely every common zero of the parent admits arbitrarily large
power-of-two choices of $b$ with an encoding $0<g<bB^W$ satisfying the
central-binomial divisibility. The corresponding signed-source interface
holds on the bounds required by the retained Sun bridge.
\end{lemma}
\begin{proof}
Every coefficient of $\mathcal D$ is a single signed weighted coefficient,
by Lemma~\ref{lem:isolation}; overlapping rows introduce no coefficient
sum. Its leading coefficient is positive. If $B>M_*+1$, the total
absolute contribution of all lower positions at $B$ is smaller than one
unit in the leading position, so $\mathcal D(B)>0$.

For $0\le x_i<b$, all coefficients of $\mathcal C$ are nonnegative and
their sum is $(1+\sum_i x_i)^2\le(32b)^2$. Therefore
\[
 \left|[T^v]\mathcal C(T)\mathcal D(T)\right|
 \le M_*(32b)^2<B/4.
\]
Collisions among absent monomials do not change this sum bound. Adding
$B/2$ at every position through $K_*$ yields digits strictly between
$B/4$ and $3B/4$, hence between zero and $B$.

The elementary block estimates retained from the source proof give
$0\le\mathcal S,\mathcal T<N$ on the soundness range $g<4bB^W$.
For example $1+g\le4bB^W$, $g<N_0$, and
\[
 2(1+g)^2\mathcal D(B)
 <32M_*b^2\frac{B^{K_*+1}}{B-1};
\]
the choice of $\Lambda$ makes this smaller than the allowed high block.
The geometric part of $\mathcal J$ is smaller than $2B^{K_*+1}$,
so $\mathcal J<N_1$. Also $T_*<B^{H_*+1}$, giving $\mathcal T<N$.
These estimates use $W,H_*,K_*$, not equal target spacing.

For $N=2^r$ and $0\le S,T<N$, the identity
\[
 (N-1)((S+T+1)N+T+1)
 =(S+T)N^2+(N-1-S)N+(N-1-T)
\]
gives
$\pc(R)=2r-\tau_2(S,T)$, where $\tau_2$ is the number of binary carries
with multiplicity. Legendre's factorial formula gives
$v_2\binom{2R}{R}=\pc(R)$. Thus the divisibility is equivalent to no
binary carries between $\mathcal S$ and $\mathcal T$.

The low-block mask first forces
$g=\sum_i x_iB^{w_i}$ with $0\le x_i<b$. The high-block mask can be
divided by two together with $\mathcal J$, which is a common binary
shift. At each selected position it tests all lower bits of the shifted
digit. Within $(B/4,3B/4)$ the only multiple of $B/2$ is $B/2$ itself.
By Lemma~\ref{lem:isolation}, this is equivalent to the vanishing of
every residual. Other positions need no additional condition.

Conversely choose $b$ larger than all coordinates of a common parent
zero. Its encoding is positive because the all-zero tuple is not a parent
zero; it is smaller than $bB^W$. Both masks are then disjoint from their
encoded blocks. The signed source chooses arbitrarily large even powers
of two $b=1+3f$ and applies the same bounds. Its congruence requirement
$3\mid R$ follows from $B\equiv N\equiv1\pmod3$, since the exponent
2140030 is even. These are the source conditions used by the inherited
ten-integer arithmetic theorem.
\end{proof}

For arbitrary nine-variable outer assignments, positivity of $\mathcal D(B)$
and the masks also gives $R\ge N\ge8$ and $0<b\le R,N$. These hold
before a putative solution is assumed. This order matters when invoking
the arithmetic bridge to establish the eventual bound on $g$.

\section{The nine-variable equation}
\label{sec:nine}
Initially use mixed domains
\[
 \epsilon,\phi,n\in\N,\qquad g,h,s,w,i,j\in\Pos.
\]
There are exactly nine unknowns. Using the source $R,N,b,g,G_*$, define
the unchanged Jones registers
\begin{align*}
 U&=N^2w, &Y&=N^2s, &M&=RY,\\
 P&=2M^2U, &K&=R+1+h(P-1), &A&=M(U+1),\\
 B_1&=2R+1, &W_1&=bw, &C&=B_1+W_1+\phi,\\
 J&=4A-5, &Q&=3W_1C-2(W_1^2-1),\\
 D&=(A^2-1)C^2+1, &e&=iJC, &E&=CDe,\\
 F&=(A^2-1)E^2+1, &G&=A+F(CD+1-A),\\
 H&=B_1+jC, &I&=(G^2-1)H^2+1,\\
 a_0&=(P^2-1)K^2+1, &a_1&=DFI,\\
 u&=JF, &v&=QF+J(H-C),\\
 m&=\bigl(K^2-4(C-KY)^2\bigr)G_*-8gK^2,
 &M_d&=Q+H.
\end{align*}
The quantity $m$ is a polynomial register, not another witness.
The numerical Jones ratio theorem says that the existential conditions
\begin{equation}\label{eq:ratio_conditions}
 a_0,a_1\in\Sq(\Z),\qquad u\mid v,\qquad m>0
\end{equation}
recover the central-binomial divisibility, a power-of-two base, and
$8g<G_*$. Conversely the source's stronger $32g<G_*$ supplies these
conditions with appropriate arithmetic witnesses. We retain the precise
\texttt{ratio\_sound} and \texttt{ratio\_complete} theorems, not an
unmodified theorem with the older exponential source positions.

\subsection{A smaller separating multiplier}
Put $B_*=P(P+R)$. The source bounds imply $U\ge64$, $M\ge512$, and
$A\ge33280$. In particular all the following inequalities hold for
\emph{every} assignment in the specified domain, before imposing any
square or divisibility condition:
\begin{equation}\label{eq:rootbounds}
 0<\sqrt{a_0}<PK\le hB_*<h\sqrt{a_1}.
\end{equation}
Indeed, $h(P+R)-K=(h-1)(R+1)\ge0$, and
$a_0=P^2K^2-K^2+1<P^2K^2$. Also
\[
 A^2-2P=M^2(U-1)^2>0,\qquad D>A^2>2P,\qquad E>D.
\]
It follows that $F>2P^2>B_*$, since $R<P$.
The quantity $L=CD+1-A$ is greater than one, so $G>F>B_*$.
Finally $I\ge G^2$ and $DF>1$, proving the last bound in
\eqref{eq:rootbounds}. Thus the existing coordinate $h$ already separates
the two radicals. No new unknown is needed for that purpose.

\subsection{The principal polynomial majorant}
Define $\lambda=2048$ and the following polynomial registers:
\begin{align*}
 L&=CD+1-A,& \mathcal N&=2G+L,\\
 \mathcal N_1&=2e^2D(D-1)L+3C,&
 \mathcal N_2&=2e^2(D-1)L,\\
 \mathcal N_3&=2e^2(L-C),& \mathcal N_4&=2e^2C,\\
 c_*&=256i(e^2+2)+2048,& B_g&=2P(P+R)+1.
\end{align*}
Put
\begin{align}
 T={}&e\bigl((1024D^2-512D-128)\mathcal N
       -64\mathcal N_1-40\mathcal N_2
       -28\mathcal N_3-21\mathcal N_4\bigr)+c_*,\label{eq:Tmajor}\\
 V={}&TH+B_g,\qquad b_1=\lambda^2a_1.\label{eq:Vmajor}
\end{align}
These definitions have no division. The registers $\mathcal N_k$ are
explicit products and sums, not quotient operations.

\begin{lemma}[Uniform polynomial majorant]\label{lem:majorant}
For every allowed assignment,
\[
 T>\lambda\sqrt{DF}\,G>0,
 \qquad \sqrt{b_1}<V-B_*,
 \qquad B_*<\sqrt{b_1}.
\]
\end{lemma}
\begin{proof}
It is enough for the first inequality that $A\ge3$ and $C,i\ge1$.
Let $X_*=e^2D^3(D-1)$. Direct substitution gives
\[
 DF=X_*+D,\qquad G=(X_*/D)L+CD+1.
\]
The fractions in this proof are real-number comparisons only.
Define the two integer polynomials
\begin{align*}
 t_6(D)={}&1024D^6-512D^5-128D^4-64D^3-40D^2-28D-21,\\
 d_5(D)={}&33792D^5+10560D^4+4928D^3+2464D^2+1176D+441.
\end{align*}
Their exact identity is
\begin{equation}\label{eq:tnidentity}
 t_6(D)^2-1024^2D^{11}(D-1)=d_5(D).
\end{equation}
For $D\ge1$, $t_6(D)>0$: the sum of the negative coefficients is 793,
less than 1024, and every corresponding power is at most $D^6$.
Also $d_5(D)>0$. Hence, for $D>1$,
\[
 \frac{t_6(D)}{D^4}>1024D^2\sqrt{1-1/D}.
\]
This finite identity replaces any need to assume convergence of a series.

The exact remainders $\mathcal N=D^k\mathcal N_k+r_k$ are
\begin{align*}
 r_1&=3-A, &r_2&=3CD+3-A,\\
 r_3&=2e^2(A-1)D^2+r_2, &
 r_4&=2e^2(1-A-C)D^3+r_3.
\end{align*}
They give the identity
\[
 T=e\frac{t_6(D)}{D^4}\mathcal N+\mathcal E+c_*,
 \qquad
 \mathcal E=e\left(\frac{64r_1}{D}+\frac{40r_2}{D^2}
                    +\frac{28r_3}{D^3}+\frac{21r_4}{D^4}\right).
\]
Since $D\ge A^2\ge9$ and $e=i(4A-5)C$, elementary bounds give
\[
 \frac{eA}{D}<\frac92 i,\qquad
 \frac{eC}{D}<\frac32 i,\qquad
 \frac{e(A+C-1)}{D}<6i.
\]
For example, the first follows from
$eA/D<4iA^2/((A^2-1)C)\le9i/2$; the second follows from
$eC/D<4iA/(A^2-1)\le3i/2$. Both $r_2$ and $r_3$ are positive.
Bounding only the possibly negative terms therefore yields
\begin{align*}
 -\mathcal E
 &<\frac{64eA}{D}+\frac{42e^3(A+C-1)}{D}+\frac{21eA}{D^4}\\
 &<252ie^2+\frac{765}{2}i<c_*-2048.
\end{align*}

For positive $X_*,D$, strict concavity of the square root gives
$\sqrt{X_*+D}<\sqrt{X_*}+D/(2\sqrt{X_*})$.
Multiplication by $2048G$ and the displayed identity for $G$ give
\[
 \lambda\sqrt{DF}\,G
 <1024eD^2\sqrt{1-1/D}\,(2G+L)+\lambda v_*,
 \quad
 v_*:=\frac{CD+1}{2eD\sqrt{1-1/D}}<1.
\]
The last inequality follows from $e\ge7C$ and
$\sqrt{D(D-1)}>D-1$. Combining these bounds proves
$T>\lambda\sqrt{DF}\,G$.
Since $I\le G^2H^2$, equations \eqref{eq:Vmajor} imply
$V-\sqrt{b_1}>B_g>B_*$. The remaining lower bound follows from
\eqref{eq:rootbounds}.
\end{proof}

The low-degree padding $B_g=2B_*+1$ has two roles. It supplies the strict
root gap for all inputs, and its factor two prevents the highest terms of
the final norm from cancelling. Omitting that factor preserves the gap
but invalidates the stated exact-degree calculation.

\subsection{The final equation and its exact existential meaning}
Set
\begin{align*}
 \Omega&=um,\\
 \Theta&=u(n+M_d)+v-\Omega(2M_d+hV).
\end{align*}
The principal nine-variable equation is
\begin{equation}\label{eq:ninefinal}
 \boxed{F_9=
 \bigl(\Theta^2-\Omega^2(a_0+h^2b_1)\bigr)^2
       -4\Omega^4h^2a_0b_1=0.}
\end{equation}
Together with the displayed source, register definitions, fixed coefficient
bytes, weights, and targets, this is a completely specified ordinary
integer polynomial in exactly nine unknowns.

\begin{proposition}[No lost sign or margin condition]\label{prop:ninenorm}
For fixed values of the first eight mixed-domain coordinates,
\eqref{eq:ninefinal} admits $n\in\N$ if and only if
$a_0,a_1$ are integer squares, $u\mid v$, and the computed margin $m$ is positive.
\end{proposition}
\begin{proof}
Here $u>0$ and $0<v/u<M_d-1$. The latter follows directly from
$v/u=Q/J+(H-C)/F$, with $Q>0$, $J,F>1$, and $H-C>0$.
If $m=0$, the left-hand polynomial is
$(u(n+M_d)+v)^4>0$. Thus a zero has $\Omega\ne0$.
Factoring the four-root norm over the real numbers gives
\[
 \frac\Theta\Omega=\varepsilon\sqrt{a_0}
                       +\delta h\sqrt{b_1},
 \qquad \varepsilon,\delta\in\{-1,1\}.
\]
This rational sum is nonzero by \eqref{eq:rootbounds}. If it is $r$, then
\[
 \sqrt{a_0}=\frac{r^2+a_0-h^2b_1}{2\varepsilon r}.
\]
Both roots are therefore rational and hence integral. This handles equal
or dependent square classes as well. Because $b_1=\lambda^2a_1$,
integrality of $\sqrt{b_1}$ implies integrality of $\sqrt{a_1}$: a rational
square root of an integer is an integer.

Write the signed integer roots as $r_0,r_1$ and define
\[
 Z=2M_d+hV+r_0+hr_1.
\]
The strict root bounds, with integer endpoints, give
$r_0\ge-hB_*+1$ and $r_1\ge-V+B_*+1$. Consequently
$Z\ge2M_d+h+1>0$. Rearranging the equation for $\Theta$ gives
\[
 n+M_d+v/u=mZ.
\]
Its left side is positive, so $m>0$. Its right side and $n+M_d$ are
integers, so $u\mid v$.

Conversely, suppose the three conditions hold. Choose either signs of
the integral roots, retain the \emph{actual} positive computed value of $m$,
and set $n=mZ-M_d-v/u$. The same bounds imply $n\ge0$, and substitution
makes \eqref{eq:ninefinal} vanish. No step replaces the computed margin
by a freely chosen witness or assumes $m=1$.
\end{proof}

Composing Proposition~\ref{prop:ninenorm} with the numerical ratio bridge
and the source lemma proves the nine-variable existence equivalence.

\subsection{A smaller-cost alternative}
The release also supplies the simpler choice $\lambda=8$,
\[
 S_8=e(8D^2-4D-1),\qquad V_8=S_8GH,\qquad b_1=64a_1.
\]
Use the same norm \eqref{eq:ninefinal} with $V=V_8$.
Its exact defect is
\[
 V_8^2-64a_1
 =\bigl(e^2(8D+1)-64D\bigr)G^2H^2+64DF(H^2-1).
\]
Since $e\ge7$, the first bracket is greater than $4De^2$.
Also $V_8<8D^2eGH$ and $2B_*<D^2$, so
\[
 \frac{V_8^2-64a_1}{V_8}
 >\frac{eGH}{2D}>\frac{(D-1)De^3}{2}>D^2>2B_*.
\]
As $V_8>B_*$, this implies $(V_8-B_*)^2>64a_1$ and proves the same
strict gap. The preceding norm proof applies unchanged.
This alternative uses 568 operations with uniform targets or 624 with
overlapping targets; its exact degrees are respectively 24,399,300 and
18,262,476. It is included because operation count and degree have
different minima among the supplied constructions.


For nine strictly positive coordinates, replace $\epsilon,\phi,n$ by
their positive coordinates minus one and leave the other six unchanged.
For nine nonnegative coordinates, leave those three unchanged and replace
$g,h,s,w,i,j$ by their nonnegative coordinates plus one. These are exact
domain bijections. The emitted circuits charge all shifts and use the
same fixed coefficient tuple.

\section{The ten-unrestricted-integer equation}
\label{sec:ten}
This branch uses $(f,g,h,k,l,w,x,y,z,m)\in\Z^{10}$ and the signed
source $b=1+3f$, $B=\Lambda b^2$, $X=R$, $Y=N^2$, $C_*=bB^W$.
The following registers are retained from the fully explicit package-40
Sun adapter:
\begin{align*}
 L&=lY, &U&=2LX, &V&=4gwY, &W_0&=bw,\\
 A&=U(V+1), &B_1&=2X+1, &C&=B_1+(A-2)h,\\
 P&=U^2V, &K&=X+1+k(P-2),\\
 S_0&=2A-5, &T_0&=3W_0C-2(W_0^2-1).
\end{align*}
\begin{align*}
 r_1&=(P^2-4)K^2+4,\\
 \delta&=8C_*^3gK^2-g^2\bigl(32(C-KL)^2C_*^3+g^2K^2\bigr).
\end{align*}
\begin{align*}
 O&=f^2l^2\delta, &\xi&=zx,\\
 D&=(A^2-4)C^2+4, &E&=C^2D\xi,\\
 F&=4(A^2-4)E^2+1,\\
 G_0&=2(A+2)(A-2)^2E^2-1, &\Omega_0&=16E,\\
 H_0&=C+(B_1-C)F+\Omega_0Fy,\\
 I_0&=(G_0^2-1)H_0^2+1.
\end{align*}
Put
\begin{align*}
 \Sigma&=\xi S_0^5,& r&=T_0+S_0y,\\
 \vartheta&=(r^3-rS_0^2-S_0^3)(r^2+23S_0^2),\\
 r_0&=b\bigl(1+(3O-4)b^2z^2\bigr),&r_2&=DFI_0,\\
 w_0&=1,&w_1&=r_0,&w_2&=r_0(P^2+K^2+3),\\
 Z&=\Sigma m-\vartheta,&A_i&=\Sigma^2w_i^2r_i\quad(0\le i\le2),\\
 V_*&=Z^2+A_2-A_0-A_1,&
 W_*&=V_*^2+4Z^2A_2-4A_0A_1.
\end{align*}
The equation is
\begin{equation}\label{eq:ten}
 \boxed{F_{10}=W_*^2-16V_*^2Z^2A_2=0.}
\end{equation}
No square root, division, binomial coefficient, or variable exponent occurs
in its definition. Those objects occur only in the correctness proof.

\subsection{What is reused in this branch}
The retained finalizer is the product over eight sign choices of a
three-radical expression, expanded into the displayed norm. Its proof
recovers the square predicates and merged divisibility condition using
separating weights, including dependent square classes. The case
$\Sigma=0$ yields $\vartheta^8\ne0$ and is excluded in the inherited
argument. The polynomial $(t^3-t-1)(t^2+23)$ supplies the cited
divisibility--nonvanishing merger, with full residue interpolation from
Sun's theorem \cite{sun,package40}.

We replace only its source. Section~\ref{sec:source} verifies the source's
positivity, range, base, congruence, and carry requirements for the new
positions. The finalizer first establishes the hypotheses needed by that
source; a square base alone is never treated as a power of two.
The inherited signed-tail theorem then proves that \eqref{eq:ten} has an
integer zero exactly when the quadratic parent has a positive zero.
This branch's new contribution is the certified sparse and overlapping
source, not a new proof of the general ten-integer theorem.

\section{Why the composition is a ZFC consistency equivalence}
\label{sec:foundations}
The complete mathematical chain is
\begin{align*}
 \text{new polynomial has a zero}
 &\Longleftrightarrow \text{the encoded quadratic system has a zero}\\
 &\Longleftrightarrow \text{the literal K215 machine halts on blank tape}\\
 &\Longleftrightarrow \text{ZFC has a finite syntactic refutation}.
\end{align*}
The first equivalence is the new encoding proof composed with the explicitly
retained numerical bridge. The latter two equivalences are the inherited
native-window/Pell compiler and the repository's calibrated proof enumerator.
They are substantive mathematical dependencies, not conclusions of running
the finite checkers in this archive.

The new arguments concern finite polynomials, integer arithmetic, finite
coefficient lists, and elementary algebraic-number reasoning. They can be
formalized in ZFC. Combining them with the identified inherited proofs
therefore supplies the asserted ZFC-provable equivalences. No assertion of
the general soundness of Lean or Rocq inside ZFC is used as a replacement
for this proof chain.

\subsection{A subtle inherited rank lemma checked in this review}
The parent uses
$I=ic^2$, $I^2=\Delta(f^2-1)$, rather than the stronger-looking expression
with an additional $\Delta$ in $I$. This must not be justified by simply
cancelling $\Delta$, because $\Delta$ need not be squarefree.
The inherited relaxed-rank lemma handles it correctly. In the notation
$\Delta=A^2-1$, $c=\psi_A(p)>A\Delta^2$, it establishes
\[
 f=\chi_A(m),\quad p\mid m,\quad c\mid m,\quad
 ic^2=\Delta\psi_A(m).
\]
Its proof passes to the fundamental unit of the squarefree part of
$\Delta$, applies strong divisibility, and uses the independently obtained
large-index bound to force the correct index divisibility. The full audit
and source locator are included in \file{proofs/FOUNDATION_AUDIT.md}.
Thus this parent does not silently assume $\Delta$ squarefree.

\subsection{Formalization status}
At the reviewed commit, the repository has an unconditional external Lean
theorem for a nine-natural-variable polynomial selected via
\texttt{Classical.choose}, and a 58-natural-variable quartic. Those are
different objects from the materialized coefficient circuits here.
The active work log records an executable coefficient compiler and an
object-language ZFC derivation as open formalization tasks \cite{status}.
The short explicit P67/P69/P73/P75 families have local Lean algebraic
certificates, with their end-to-end native compiler bridge still an ordinary
mathematical proof. This contribution adds ordinary proofs and independently
checked finite data; it does not claim to close those formalization tasks.

No consistency assumption is used to define the coefficients. The theorem
proves an equivalence and proves neither side. A claim of two-sided
independence would require additional hypotheses and is not made here.

\section{Reproduction and the scope of the evidence}
\label{sec:verification}
The release contains Python code using only the standard library. The
recommended entry point is \file{verify.py}; \file{verify_refined.py}
checks the refined nine-variable circuits. The independent support checker
does not import construction helpers. The scripts are also run with
\texttt{python3 -O}, so correctness gates must not disappear when assertions
are disabled. The archive includes exact receipts and a complete hash manifest.

The checks establish the following finite claims:
\begin{enumerate}
\item The parent expansion and all successor shifts match the fixed
31-coordinate, 22-equation source.
\item All 71 present monomials are isolated against all 528 possible
quadratic monomials, and every cross-row target is isolated in both layouts.
\item The four compiler ports and all derived coefficient recipes are
bound to the supplied bytes; source positivity and coefficient bounds use
that complete tuple.
\item The emitted straight-line graphs are acyclic, every gate and
witness port is live, and every operation is charged under the stated model.
\item Direct source-formula replay matches the emitted graphs; exact
symbolic identities relate the finalizer to its cancellation-resolved
degree certificate.
\item Multivariate degree upper bounds are attained, by a nonzero
diagonal leading coefficient, with exact formulas for the nine-variable
tail and independent modular certificates for the ten-variable tail.
\end{enumerate}
These are checks of the polynomial construction and its finite certificates.
Searching sample assignments cannot prove an unbounded consistency
equivalence. The universal claims instead follow from the mathematical
proofs above and their named inherited dependencies.

\section{Where the reduction comes from and what remains}
The first gain comes from specializing the algebraic encoding to the
\emph{support} of the actual residual system. The 457 absent quadratic
monomials never need to be separated from one another. This is stronger
than merely searching for a better full Sidon set, because it changes
the constraint being solved.

The second gain separates only the output positions that will actually
be tested. The row-polynomial supports may overlap freely elsewhere.
The explicit target list is a small exact certificate for this weaker
condition, while every carry bound is preserved.

The third gain uses divisibility already present in the arithmetic tail
to approximate its radical more closely by a polynomial. It also replaces
a large separating multiplier by the existing witness $h$, after proving
an unconditional gap. These transformations introduce no witness and no
variable denominator. Their defect identities explain the lower degree
and provide small exact certificates of the otherwise hidden cancellations.

Further work can optimize the joint choice of variable weights, row order,
and charged power schedule; formalize Lemma~\ref{lem:isolation} once for
all finite sparse sources; and connect executable coefficient data to the
existing Lean ZFC endpoint. None of those possible improvements is included
in the certified numerical claims of this release.

\appendix
\section{Exact finite placement data}
The overlap table gives the original parent row number, not its index in the reordered construction array.
\begin{center}
\begin{tabular}{@{}rrrr@{}}
\toprule Target index $j$ & $H_j$ & Original row & Sign\\\midrule
0 & 358 & 14 & $+1$\\
1 & 608 & 7 & $+1$\\
2 & 939 & 21 & $+1$\\
3 & 1147 & 10 & $+1$\\
4 & 1432 & 8 & $+1$\\
5 & 1739 & 17 & $+1$\\
6 & 1989 & 16 & $+1$\\
7 & 2170 & 6 & $+1$\\
8 & 2512 & 1 & $+1$\\
9 & 2699 & 13 & $+1$\\
10 & 2972 & 4 & $+1$\\
11 & 3305 & 2 & $+1$\\
12 & 3492 & 19 & $+1$\\
13 & 3683 & 11 & $+1$\\
14 & 3914 & 20 & $+1$\\
15 & 4108 & 5 & $+1$\\
16 & 4399 & 18 & $+1$\\
17 & 4697 & 15 & $+1$\\
18 & 4926 & 22 & $+1$\\
19 & 5224 & 3 & $+1$\\
20 & 5469 & 12 & $+1$\\
21 & 5775 & 9 & $-1$\\
\bottomrule
\end{tabular}
\end{center}
The uniform placement uses all positive signs and original row order $1,2,3,4,6,7,\ldots,22,5$, with $H_j=358+359j$. Both layouts use the weights in equation~\eqref{eq:weights}.

\section{Detailed degree certificate}
Let \(t=\deg N=3(H_*+3W+2)\). Thus \(t=25\,308\) for uniform targets
and \(t=18\,942\) for overlapping targets. All degrees below are
multivariate total degrees. A \emph{leader} is the leading coefficient
after substituting the same indeterminate for all nine coordinates.
A nonzero such coefficient proves that the multivariate upper bound
is attained. Successor/predecessor domain shifts leave these leaders unchanged.

\begin{center}
\begin{tabular}{@{}lll@{}}
\toprule Register & Degree & Diagonal leader\\\midrule
\(N\) & \(t\) & \(2^2\)\\
\(R\) & \(3t\) & \(2^4\)\\
\(P\) & \(12t+3\) & \(2^{21}\)\\
\(A\) & \(7t+2\) & \(2^{12}\)\\
\(C\) & \(3t\) & \(2^5\)\\
\(D\) & \(20t+4\) & \(2^{34}\)\\
\(e\) & \(10t+3\) & \(2^{19}\)\\
\(L\) & \(23t+4\) & \(2^{39}\)\\
\(F\) & \(80t+18\) & \(2^{140}\)\\
\(G\) & \(103t+22\) & \(2^{179}\)\\
\(H\) & \(3t+1\) & \(2^5\)\\
\(a_0\) & \(48t+14\) & \(2^{84}\)\\
\(a_1=DFI\) & \(312t+68\) & \(2^{542}\)\\
\(\Omega\) & \(115t+3W+31\) & \(-2^{211}\)\\
\(T\), principal majorant & \(153t+33\) & \(2^{277}\)\\
\(V\), principal majorant & \(156t+34\) & \(2^{282}\)\\
\bottomrule
\end{tabular}
\end{center}

\subsection{A polynomial certificate for the hidden cancellation}
Use \(t_6,d_5,r_k,\mathcal N,L,c_*\) from the majorant proof and define
\[
 R_*=
 e(64r_1D^3+40r_2D^2+28r_3D+21r_4)+c_*D^4 .
\]
Then \(D^4T=e\,t_6\mathcal N+R_*\), and the following are exact
integer-polynomial identities:
\begin{align}
 \Delta_T:={}&D^8(T^2-\lambda^2DFG^2)\notag\\
 ={}&e^2d_5\mathcal N^2
 +1024^2D^8\bigl(X_*L^2-4DG(CD+1)\bigr)\notag\\
 &\quad+2e\,t_6\mathcal N R_*+R_*^2,\label{eq:deltaT}\\
 \Delta_V:={}&D^8(V^2-\lambda^2a_1)\notag\\
 ={}&H^2\Delta_T+\lambda^2DFD^8(H^2-1)
              +D^8(2THB_g+B_g^2).\label{eq:deltaV}
\end{align}
Here \(X_*=e^2D^3(D-1)\). The term
\(X_*\mathcal N^2-4(X_*+D)G^2\) simplifies to
\(X_*L^2-4DG(CD+1)\); this is the essential cancellation
in \eqref{eq:deltaT}.

The respective degree bounds of the four summands in
\eqref{eq:deltaT} are
\[
326t+70,\quad306t+62,\quad333t+72,\quad200t+46.
\]
For the third bound, \(\deg R_*=100t+23\), with a positive leading
coefficient supplied by \(c_*D^4\). Its term uniquely dominates.
Since \(\deg D^8=160t+32\), we have
\[
\deg(T^2-\lambda^2DFG^2)=173t+40.
\]
After multiplication by \(H^2\), that becomes \(179t+42\).
In \(V^2-\lambda^2a_1\), however, the padding cross term \(2THB_g\)
has degree \(180t+40\), which is greater because \(t>2\).
Its leader is
\(2\cdot2^{282}\cdot2^{43}=2^{326}\).
All other terms have smaller degree. Thus
\[
 \deg(V^2-\lambda^2a_1)=180t+40,\qquad
 \operatorname{lc}_{\rm diag}(V^2-\lambda^2a_1)=2^{326}.
\]

Let \(E_0=u(n+M_d)+v-2\Omega M_d\). The inner norm has the exact expansion
\[
 \Theta^2-\Omega^2(a_0+h^2b_1)
 =E_0^2-2E_0\Omega hV
       +\Omega^2h^2(V^2-b_1)-\Omega^2a_0.
\]
The third term uniquely dominates: its degree is \(410t+6W+104\)
and its diagonal leader is \(2^{748}\). Squaring produces degree
\(820t+12W+208\) and leader \(2^{1496}\).
The subtracted norm term has the \emph{same} degree and leader
\(2^{1494}\). Therefore
\[
 \boxed{\deg F_9=820t+12W+208,\qquad
        \operatorname{lc}_{\rm diag}F_9=3\cdot2^{1494}.}
\]
This proves exact degree, not merely a syntactic upper bound.
It also explains the doubled padding \(B_g\): padding only by \(B_*+1\)
would make the two highest terms cancel.

\subsection{What the certificate graph evaluates}
The primary graph evaluates exactly \(F_9\). The principal degree
certificate graph instead evaluates \(D^{16}F_9\), through
\[
 \left[
 D^8(E_0^2-2E_0\Omega hV-\Omega^2a_0)
 +\Omega^2h^2\Delta_V
 \right]^2
 -4D^{16}\Omega^4h^2a_0b_1.
\]
This avoids division while exposing every required cancellation.
Its exact degree is \(1140t+12W+272\) and its diagonal leader is
\(3\cdot2^{2038}\). Subtracting the known degree of \(D^{16}\)
recovers the primary degree. The certificate graph is not substituted
for the delivered equation.

The V8 graph uses an equal-polynomial certificate instead. Its displayed
defect has degree \(252t+56\) and diagonal leader \(2^{443}\).
Consequently
\[
 \deg F_{9,\mathrm{V8}}=964t+12W+240,\qquad
 \operatorname{lc}_{\rm diag}F_{9,\mathrm{V8}}=2^{1730}.
\]

\subsection{Unchanged-finalizer comparison and modular receipts}
For the source replacements with the inherited finalizers, the degree
formulas are
\begin{align*}
 \deg F_{9,\mathrm{inherited}}&=3456H_*+10380W+7192,\\
 \deg F_{10}&=6320H_*+19016W+12992.
\end{align*}
The inherited-finalizer graphs and their nonzero leading coefficients are
independently replayed modulo \(1\,000\,000\,007\), \(1\,000\,000\,009\), and
\(998\,244\,353\). The refined graphs use \(1\,000\,003\),
\(1\,000\,033\), and \(1\,000\,000\,007\).
A single nonzero residue at a proved multivariate upper bound suffices
for exactness; the three receipts provide additional cross-checks.
For the principal nine-variable graph the symbolic degree calculation
above also supplies an exact integer leader.


\section{Files, fixed coefficients, and provenance}
The machine-readable equations live in \file{generated/}. The primary
refined nine-variable equations and their degree
certificates live in \file{refined/}; the principal certificate evaluates
$D^{16}F_9$, while the V8 certificate evaluates $F_9$ itself.
\file{data/} contains the literal
parent, complete ports, coefficient recipes, weights, and placement data.
\file{proofs/} contains the independent mathematical reviews; \file{upstream/}
and \file{inherited/} contain the pinned dependency text needed to inspect
the inherited interfaces.
The source and dependency hashes are recorded in the manifests. The compiler
ports are immutable data, never additional unknowns.

The manuscript source is \file{report.tex}; a normal TeX installation can
rebuild it with two passes of \texttt{pdflatex}. The adjacent small TeX input
files contain the exact tables and the finalizer proof. They are part of the
source deliverable. New material is supplied under the repository's MIT
No Attribution license; inherited notices and attribution remain in the archive.

\begin{thebibliography}{9}
\raggedright
\bibitem{repo} V. Reshetnikov and ProveIt contributors.
\emph{ProveIt}, reviewed commit
\texttt{124c07c0b77a9e20e84676514d8b1ff78536ea3f}, 23 October 2026.
\href{https://github.com/VladimirReshetnikov/ProveIt/tree/124c07c0b77a9e20e84676514d8b1ff78536ea3f/Computability/DiophantineConsistency}{Pinned DiophantineConsistency repository directory}.
\bibitem{package40} ProveIt report package 40.
\emph{Direct quadratic packing for a ten-integer consistency equation},
23 October 2026. Pinned repository directory
\texttt{con-zfc-diophantine-forms/40-packing10}; report and source certificates
at the commit in \cite{repo}.
\bibitem{package45} ProveIt report package 45.
\emph{Simultaneous quadratic coding for ZFC consistency}, 23 October 2026.
Pinned directory \texttt{con-zfc-diophantine-forms/45-multioutput};
especially the source proof and \texttt{MULTITARGET\_NINE\_PROOF.md}.
\bibitem{parent} ProveIt report package 31.
\emph{Sun ten-integer and quadratic31 constructions}; literal parent,
full coefficient tuple, and retained-system proofs in
\texttt{con-zfc-diophantine-forms/31-sun10}, at \cite{repo}.
\bibitem{status} ProveIt contributors. \emph{DiophantineConsistency research
status}, \texttt{STATUS.md}, and \emph{Constructing the coefficient data},
\texttt{COMPILER.md}, at \cite{repo}.
\bibitem{jones} J. P. Jones. Universal Diophantine equation.
\emph{Journal of Symbolic Logic} \textbf{47}(3) (1982), 549--571.
\url{https://doi.org/10.2307/2273588}.
The numerical ratio interfaces used here are also formalized in ProveIt's
\texttt{Diophantine/Paper1982/NineRatio.lean}.
\bibitem{sun} Z.-W. Sun. \emph{Ten unknowns for Hilbert's tenth problem
over the integers}. arXiv:2609.01594v2, 2 September 2026.
\url{https://arxiv.org/html/2609.01594v2}.
\bibitem{relaxed} ProveIt contributors. \emph{Relaxed Pell auxiliary proof}.
\nolinkurl{Computability/HilbertTenthProblem/Papers/1980/PELL_RELAXED_AUXILIARY_PROOF.md},
at \cite{repo}; specialized in package 36's inherited ordinary proof.
\end{thebibliography}
\end{document}
